\documentclass[18]{article}
\usepackage{amsmath}
\usepackage{amsthm}

\begin{document}
\begin{center}
\Large\textbf{Assignment 2-Math}
\end{center}
\begin{center}
\large{TW Lab 2025}
\end{center}
\vspace{1cm}e
\begin{center}
Write \LaTeX code to output the following sections as it is .Use \LaTeX reference wherever required. You may also explore the Math menu provided as Menu bar
\end{center}
\section{Basics}
This famous Einstein's mass-energy equation,$E=mc^2$ should be rendered inline.However the following equation of motion should be rendered in a block:
\[
v = u + at
\]

\section{Fraction \& Roots}
\[
x_i={\frac{\sqrt{a^2+b_i^2}}{c_i}}
\]

\section{Symbols \& Math-fonts}
\[
G_{\mu\nu}+\Lambda g_{\mu\nu}=\kappa T_{\mu\nu},   \forall x \in \mathbf{R} \ \exists \varepsilon > 0 \Rightarrow |x| < \varepsilon
\]
\[
\mathcal{F}(\theta)=\sum_{k=0}^{n}{n \choose k} \theta^k (1-\theta)^{n-k}
\]

\section{Calculus}
\[
\frac{du}{dx}+\frac{\partial u}{\partial t}=\alpha \frac{\partial^2 u}{\partial^2 x}
\]
\[
\int_{0}^{1}(3x^2 + 2x+1)dx
\]
\vspace{4cm}
\section{Multi-line equations}
\subsection{Centering \& alignment}
\begin{equation}
x=5
\end{equation}
\begin{equation}
x+y=3
\end{equation}
\vspace{1cm}
\begin{align}
x=5
\end{align}
\begin{align}
x+y=3
\end{align}
\subsection{Multiline equation}
\begin{multline*}
(x+y+z)^3=(x+y+z).(x+y+z).(x+y+z)=x^3+y^3+z^3+ \\
3x^2y+3x^2y+3y^2x+3y^2z+3z^2x+3z^2y+6xyz
\end{multline*}
\section{Matrices}
\[
\begin{matrix}
1 & 2 \\
3 & 4
\end{matrix}
\quad
\begin{bmatrix}
1 & 2 \\
3 & 4
\end{bmatrix}
\quad
\begin{vmatrix}
1 & 2 \\
3 & 4
\end{vmatrix}
\quad
\begin{pmatrix}
1 & 2 \\
3 & 4
\end{pmatrix}
\]
\end{document}
